\documentclass[12pt]{article}
\usepackage[T1]{fontenc}
\usepackage[utf8]{inputenc}
\usepackage{lmodern}
\usepackage{textcomp}
\usepackage{enumitem}
\usepackage{geometry}
\geometry{margin=1in}
\begin{document}
\begin{center}
Answer Key - Sociology - Undergraduate Level Assessment 14 \
\small (Answers are ordered exactly as in the question paper.)
\end{center}

\section*{Multiple Choice}
\begin{enumerate}[label=\arabic*.]
\item \textbf{Answer:} D
\\ \textit{Options:} A) flexibility was reduced by the introduction of detailed daily worksheets; B) the state had greater control than managers over production processes; C) ownership was transferred from small shareholders to senior managers; D) employment depended on performance rather than trust and commitment
\\ \textit{Note:} The correct answer reflects Mulholland's argument that privatization shifted the focus of employment relationships from a foundation of trust and loyalty to one based on measurable performance metrics, thereby altering the dynamics between companies and their managers.
\item \textbf{Answer:} C
\\ \textit{Options:} A) there are no clearly defined projects into which the money can be directed; B) the United Nations has refused to call on rich countries to provide it; C) debt repayments with interest can be greater than the amount of money received; D) debt repayments with interest can be greater than the amount of money received
\\ \textit{Note:} The correct answer highlights that the burden of debt repayments, often with high interest rates, can exceed the benefits of economic aid received, thereby undermining the potential for modernization in developing countries.
\item \textbf{Answer:} D
\\ \textit{Options:} A) age; B) income; C) verbal fluency; D) occupation
\\ \textit{Note:} In modern societies, occupation is a primary indicator of social status because it often determines an individual's income, prestige, and access to resources, which are critical factors in evaluating one's position within the social hierarchy.
\item \textbf{Answer:} A
\\ \textit{Options:} A) ways of acting, thinking, and feeling that are collective and social in origin; B) the way scientists construct knowledge in a social context; C) data collected about social phenomena that are proven to be correct; D) ideas and theories that have no basis in the external, physical world
\\ \textit{Note:} This answer is correct because Durkheim emphasized that social facts are phenomena that exist outside the individual and exert control over social behavior, highlighting their collective nature and social origins.
\item \textbf{Answer:} C
\\ \textit{Options:} A) cults of the capital; B) capital culture; C) cultural capital; D) culpable capture
\\ \textit{Note:} Bourdieu argued that cultural capital, which encompasses non-financial social assets such as education, tastes, and cultural knowledge, plays a crucial role in perpetuating social class distinctions and inequalities across generations.
\end{enumerate}

\section*{True / False}
\begin{enumerate}[label=\arabic*.]
\setcounter{enumi}{5}
\item \textbf{Answer:} False
\\ \textit{Options:} A) True; B) False
\\ \textit{Note:} Dahrendorf, Rex, and Habermas are primarily known for their analyses of social conflict, power structures, and the role of discourse in society rather than focusing on social solidarity and cohesion among communities.
\item \textbf{Answer:} False
\\ \textit{Options:} A) True; B) False
\\ \textit{Note:} The trend of decarceration actually involves a shift towards reducing the prison population and emphasizing alternative forms of punishment, rather than increasing imprisonment.
\item \textbf{Answer:} False
\\ \textit{Options:} A) True; B) False
\\ \textit{Note:} The social democratic model emphasizes collective welfare and social safety nets, which typically result in lower unemployment rates compared to individualistic models that prioritize self-reliance.
\item \textbf{Answer:} True
\\ \textit{Options:} A) True; B) False
\\ \textit{Note:} Fulcher and Scott argue that community can be understood as a social construct characterized by relationships and shared identities rather than being strictly tied to a specific geographical location.
\item \textbf{Answer:} False
\\ \textit{Options:} A) True; B) False
\\ \textit{Note:} A document can be considered authentic based on its origin and integrity, but it may still contain political bias or subjective interpretations, making the statement false.
\end{enumerate}

\section*{Long Form Response}
\begin{enumerate}[label=\arabic*.]
\setcounter{enumi}{10}
\item \textbf{Answer:} Scapegoating in sociology refers to the practice of blaming an individual or group for problems or misfortunes that are not necessarily their fault, often diverting attention from the actual causes. This phenomenon can have significant implications for societal dynamics, as it fosters division, reinforces stereotypes, and perpetuates social hierarchies. The targeted group often experiences marginalization and discrimination, leading to psychological and social repercussions. Additionally, scapegoating can unify the broader community around a common enemy, but it does so at the expense of justice and understanding, ultimately hindering social cohesion and progress.
\\ \textit{Note:} correct because it accurately defines scapegoating as the unjust attribution of blame, highlights its divisive effects on society and stereotypes, and addresses the negative consequences faced by both the targeted group and the wider community.
\item \textbf{Answer:} Domhoff identifies several key processes of decision-making in the USA, including the role of elites, the influence of organized interests, and the impact of institutional settings on policy outcomes. These processes emphasize how power dynamics shape decisions at various levels of governance. However, the exploitation process is notably absent from his framework, which can be critically analyzed as a significant oversight. The omission may stem from Domhoff's focus on formal decision-making structures rather than the underlying economic inequalities that exploitation represents. By excluding the exploitation process, his framework might underappreciate the ways in which socioeconomic disparities influence decision-making and policy implementation, ultimately limiting a comprehensive understanding of power relations in society.
\\ \textit{Note:} correct because it accurately outlines Domhoff's key decision-making processes, such as the influence of elites and organized interests, while also highlighting the critical gap of the exploitation process in his framework, suggesting that this oversight may be due to his emphasis on formal structures over informal power dynamics.
\end{enumerate}

\vfill
\noindent\textit{End of Answer Key}
\end{document}